\section{Qualitative topology and strong smoothing}
\label{sec:smoothing}

All PL structures in this Section are the ordinary structures on open
subsets of $\R^3$.  Triangulations are locally finite in the open domain.
In particular, neither a triangulation nor the local constants used below
need have uniform behavior at the boundary.

The qualitative PL tools belong to the three-dimensional triangulation
and approximation theory of Moise and Bing
\cite{Moise1952Approximation,Moise1952Triangulation,Bing1959}; Hamilton's
relative formulation \cite{Hamilton1976} is the precise interface used
below. Edwards--Kirby's deformation theorem supplies the supported
ambient extensions \cite{EdwardsKirby1971}. We separate these
topological existence statements from the derivative estimates:
finite local model libraries and subsequently chosen small exceptional
sets provide the latter.

\begin{theorem}[Strong smoothing of locally bi-Lipschitz maps]
\label{sm:smoothing}
Let $\Omega,\Omega'\subset\R^3$ be bounded domains, let
$1\le p<\infty$, and let $H:\Omega\to\Omega'$ be a locally
bi-Lipschitz homeomorphism in $W^{1,p}(\Omega;\R^3)$.  For every
$\eps>0$ and every continuous function $\xi:\Omega\to(0,\infty)$,
there is a $C^\infty$ diffeomorphism $g:\Omega\to\Omega'$ such that
\begin{equation}
 \label{sm:smoothing-conclusion}
 \norm{g-H}_{W^{1,p}(\Omega)}<\eps,
 \qquad |g(x)-H(x)|<\xi(x)\quad(x\in\Omega).
\end{equation}
In particular, $g$ belongs to $W^{1,p}(\Omega;\R^3)$, and the
approximations have exactly the target $\Omega'$.
\end{theorem}

The proof has two parts.  First we smooth a locally finite PL
homeomorphism with summable local derivative errors.  We then approximate
$H$ by such a PL map.  In the second part a finite collection of local
models gives derivative bounds before the measure of the exceptional
mesh cubes is made small.  The qualitative topology used to choose those
models supplies no derivative estimate.

\subsection{Relative qualitative tools}

\begin{lemma}[Fine relative PL approximation]
\label{sm:relative-pl}
Let $M,N$ be PL three-manifolds without boundary, with $N$ metrized,
and let $f:M\to N$ be a homeomorphism.  Suppose that $K\subset M$ is
closed and that $f$ is PL on an open neighborhood of $K$.  Given a
continuous function $\eta:M\to(0,\infty)$, there is a PL
homeomorphism $f_1:M\to N$ which agrees with $f$ on a neighborhood of
$K$ and satisfies
\[
 d_N(f_1(x),f(x))<\eta(x)\qquad(x\in M).
\]
Noncompact manifolds are allowed.  For manifolds with boundary the same
statement holds if $K$ contains $\partial M$ and $f$ is PL near $K$.
\end{lemma}

\begin{proof}
Choose a closed neighborhood $K^+$ of $K$ contained in the open set on
which $f$ is PL.  Transport the PL structure of $N$ to $M$ by $f$.
The identity from the original structure to the transported structure
is PL near $K^+$.  Hamilton's relative approximation Theorem
\cite[Section~3, Theorem~2.2, pp.~68--69]{Hamilton1976} applies to these
two structures.  Use the compatible metric
$d_f(x,y)=d_N(f(x),f(y))$ and the fine tolerance $\eta$.
It gives a homeomorphism $a$ that is PL between the two structures,
is the identity on $K^+$, and satisfies
$d_f(a(x),x)<\eta(x)$.  Then $f_1=f\circ a$ has all the asserted
properties.  The boundary condition in Hamilton's Theorem is precisely
the additional condition stated here; for an open Euclidean domain its
manifold boundary is empty.
\end{proof}

\begin{lemma}[Small local ambient extension]
\label{sm:local-extension}
Fix a closed Euclidean cube $B\subset\R^3$, compact sets $C,D\subset B$
with $C\subset\operatorname{int}B$, and an open neighborhood $O$ of
$C\cup D\cup\partial B$.  For every $\lambda>0$ there is a
$\delta>0$, depending only on these source sets and on $\lambda$,
with the following property.  If $e:O\to\R^3$ is an embedding,
\[
 \sup_{x\in O}|e(x)-x|<\delta,
 \qquad e=\Id\ \hbox{on }D\cup\partial B,
\]
then there is a homeomorphism $A:\R^3\to\R^3$ such that
\begin{equation}
 \label{sm:extension-conclusion}
 A=e\ \hbox{on }C,\qquad
 A=\Id\ \hbox{on }D\cup(\R^3\setminus\operatorname{int}B),
 \qquad \sup_{\R^3}|A-\Id|<\lambda.
\end{equation}
The number $\delta$ is uniform over a finite list of fixed normalized
configurations $(B,C,D,O)$.  No derivative bound on $e$ is required.
\end{lemma}

\begin{proof}
Put $E=C\cup D\cup\partial B$.  The deformation Theorem of
Edwards--Kirby \cite[Theorem~5.1]{EdwardsKirby1971}, applied at the
inclusion of $O$ in $\R^3$, deforms every sufficiently close embedding
$e$ through embeddings $e_t$, starting at $e_0=e$, so that
$e_1=\Id$ on $E$.  The deformation is unchanged outside a fixed
compact neighborhood $L$ of $E$ in $O$, and fixes the inclusion.
Its continuity at the inclusion permits us to require
\[
 |e_t(x)-x|<\lambda/2\quad(x\in L,\ 0\le t\le1),
 \qquad |e(x)-x|<\lambda/2\quad(x\in L).
\]
A sufficiently small uniform $\delta$ on $O$ meets the finitely many
compact-open conditions that specify this neighborhood of the
inclusion.  In the terminology of that Theorem a proper embedding
respects manifold boundaries; this condition is automatic for the
ambient manifold $\R^3$.

All the open images $e_t(O)$ coincide.  To see this, enclose $L$ in a
finite union of relatively compact subdomains of $O$ whose boundary collars are outside
the region where the deformation changes the map.  The embeddings
agree on that collar. Invariance of domain
\cite[Theorem~2B.3]{Hatcher2002} and degree relative to
the common boundary show that the images of the enclosed subdomain
are the same.  Outside it the maps already agree.
On this common open image define
\[
 A_0=e\circ e_1^{-1},
\]
and define $A_0$ to be the identity elsewhere.  The map is already
the identity off the compact set $e_1(L)$ in the common image, so
this is an ambient homeomorphism.  It agrees with $e$ on $E$ and has
displacement less than $\lambda$.  Since it fixes $\partial B$
pointwise, it preserves the bounded complementary component
$\operatorname{int}B$.  Restrict $A_0$ to $B$ and extend it by the
identity outside $B$ to obtain $A$.

The construction uses only the displayed source configuration and
continuity at the inclusion.  For finitely many normalized
configurations take the minimum of the resulting positive tolerances.
\end{proof}

We will use Lemma~\ref{sm:local-extension} when the input embedding is
defined by a small map $\eta$ near a new core and by the identity near
older protected sets.  The following observation specifies the required
gluing condition.  Choose open sets
$\overline{O_0}\subset U_0$, where $\eta$ is defined on $U_0$, and
an open set $V$ containing the protected sets, such that
$\eta=\Id$ on $U_0\cap V$.  If $\eta$ is sufficiently close to
the identity that $\eta(O_0)\subset U_0$, the map
\begin{equation}
 \label{sm:patched-embedding}
 e=\eta\ \hbox{on }O_0,\qquad e=\Id\ \hbox{on }V
\end{equation}
is an embedding.  Indeed, a collision $\eta(x)=y$ with
$x\in O_0$ and $y\in V$ has $y\in U_0\cap V$, whence
$\eta(y)=y$ and injectivity of $\eta$ gives $x=y$.  The open
embedding property then follows from invariance of domain.  All the
sets in this observation can be fixed in advance when the source
configurations range over a finite list.

\begin{lemma}[Fine control, properness, and the target]
\label{sm:proper-degree}
Let $h:\Omega\to\Omega'$ be an orientation-preserving homeomorphism
between bounded domains, and let $g:\Omega\to\R^3$ be a smooth
map with $\det Dg>0$.  If
\begin{equation}
 \label{sm:target-distance}
 |g(x)-h(x)|<\frac12\dist(h(x),\R^3\setminus\Omega')
 \qquad(x\in\Omega),
\end{equation}
then $g$ is a diffeomorphism of $\Omega$ onto $\Omega'$.
\end{lemma}

\begin{proof}
Write $d(y)=\dist(y,\R^3\setminus\Omega')$ and
$g_t=(1-t)h+tg$.  Condition~\eqref{sm:target-distance} puts the entire
segment $g_t(x)$ in $\Omega'$.  Since $d$ is 1-Lipschitz,
\[
 d(g_t(x))\le\frac32d(h(x)).
\]
If $K\Subset\Omega'$ and $m=\min_Kd>0$, the inverse image of $K$
under the homotopy lies in
\[
 h^{-1}\bigl(\{y\in\Omega':d(y)\ge2m/3\}\bigr)\times[0,1].
\]
The set in braces is compact because $\Omega'$ is bounded.
Thus $g_t$ is a proper homotopy.  Its endpoint $g$ has the same
degree as $h$, namely one.  A proper local diffeomorphism has
finitely many preimages of each point, and here each preimage has
positive local degree.  The degree-one identity therefore says that
every point of $\Omega'$ has exactly one preimage.  The local smooth
inverses patch to a smooth inverse on $\Omega'$.
\end{proof}

\subsection{Smoothing a locally finite PL homeomorphism}

Campbell--D'Onofrio--V\'itek proved diffeomorphic approximation of
locally finite piecewise affine homeomorphisms in dimensions three
and four, with uniform value control and derivative-error control
for both the map and its inverse
\cite[Theorem~A]{CampbellDOnofrioVitek2026}.
We give the local proof needed here, including arbitrary positive
continuous value tolerances.

\begin{proposition}[Strong PL smoothing]
\label{sm:pl-smoothing}
The conclusion of Theorem~\ref{sm:smoothing} holds if $H$ is replaced
by a locally finite PL homeomorphism
$h:\Omega\to\Omega'$ in $W^{1,p}(\Omega;\R^3)$.
\end{proposition}

\begin{proof}
An orientation-reversing Euclidean isometry in the target reduces
the proof to the orientation-preserving case.  Fix a locally finite
affine triangulation for $h$.  We prescribe summable error budgets
for its edges, vertices, and a locally finite family of compact sets
covering the remainder of the domain.  All modifications below can
also obey an arbitrary positive continuous value tolerance.

\paragraph{The open edges.}
Around each open edge choose a tube of radius $d(t)$, where
$0<t<l$ is its axial coordinate.  The tubes of distinct open edges
are disjoint, their radii are bounded by a small multiple of
$\min(t,l-t)$, and $d$ is exactly linear near both endpoints.
These choices follow from the finite geometry of each simplex star;
local finiteness makes them compatible throughout $\Omega$.
We may make $\norm{d'}_\infty$ small by decreasing the width.

In positively oriented orthogonal frames, and after translations,
the original map on this tube is
\begin{equation}
 \label{sm:edge-form}
 (t,r,\theta)\longmapsto
 \bigl(at+r b(\theta),\;r s(\theta)e_{\phi(\theta)}\bigr),
 \qquad e_\alpha=(\cos\alpha,\sin\alpha),
\end{equation}
where $a>0$, $s>0$, and the coefficients are smooth on the closed
angular sectors.  The transverse homogeneous map is injective:
otherwise two equal transverse images, taken at arbitrarily small
radius, would have their axial difference canceled by changing $t$,
contradicting injectivity of $h$ in the edge star.  Its angular
map is consequently an orientation-preserving circle
homeomorphism.  We choose a lift $\phi$ with
$\phi(\theta+2\pi)=\phi(\theta)+2\pi$.
Positive determinants on the finitely many sectors give upper and
positive lower bounds for $s$ and for the one-sided derivatives
$\phi'$.

Choose constants $c>0$ and $c'$, and for $0\le\tau\le1$ put
\[
 b_\tau=(1-\tau)b,\qquad
 s_\tau=(1-\tau)s+\tau c,\qquad
 \phi_\tau=(1-\tau)\phi+\tau(\theta+c').
\]
For fixed $\tau$, the determinant of the corresponding map in the
frame $(\partial_t,\partial_r,r^{-1}\partial_\theta)$ is
\begin{equation}
 \label{sm:edge-determinant}
 a s_\tau^2\phi_\tau'>0.
\end{equation}
On all the closed sectors and all $\tau\in[0,1]$ this has a
positive minimum.  Let $\chi$ be a smooth cutoff equal to 1 on
$(-\infty,-1]$ and to 0 on $[-1/2,\infty)$, and substitute
\[
 \tau(t,r)=\chi\!\left(\frac{\log(r/d(t))}{N}\right)
\]
in the preceding formula.  The additional derivative columns have
size at most
\begin{equation}
 \label{sm:edge-cutoff-error}
 \frac{C}{N}\bigl(1+\norm{d'}_\infty\bigr).
\end{equation}
Indeed, $|r\partial_r\tau|\le C/N$ and
$|r\partial_t\tau|\le C(r/d)|d'|/N$, while $r/d\le1$ on
the support of the modification.  The remaining factors come from
the fixed angular data.  Taking $N$ large preserves positive
determinants on every sector, including one-sided limits.
Near the axis the result is the nonsingular affine map
$(t,r,\theta)\mapsto(at,cr e_{\theta+c'})$.
Near the tube boundary it agrees with $h$.

All first derivatives are bounded by an edge-dependent constant
independent of a further common decrease of $d$.  The volume of the
tube tends to zero under that decrease.  Both its derivative error
and its value error therefore have arbitrarily small $W^{1,p}$
cost.  Choose these costs summably over the edges and call the
result $h_1$.  It is continuous and locally Lipschitz, fixes all
vertices, and has the required small total error.  Because the
tube widths are exactly linear near endpoints, $h_1-h_1(v)$ is
homogeneous of degree one on a sufficiently small ball about every
vertex $v$.

\paragraph{The faces and the punctured vertex balls.}
Away from the vertices, the closed convex hull of the nearby
almost-everywhere gradients of $h_1$ is contained in
$\GL^+(3)$ on a sufficiently small neighborhood of each point.
At a smooth point this follows by continuity.  At a nonsmooth
point the only remaining interface is one face.  If $A_-$ and
$A_+$ are its two gradient traces, continuity gives agreement on
the tangent plane, so $A_+-A_-$ has rank at most one with normal
covector to that face.  Consequently
\begin{equation}
 \label{sm:face-segment}
 \det\bigl((1-t)A_-+tA_+\bigr)
 =(1-t)\det A_-+t\det A_+>0\quad(0\le t\le1).
\end{equation}
The traces vary continuously on the two sides.  Their nearby
convex hull lies in a sufficiently small neighborhood of this
compact positive-determinant segment, proving the assertion.
Homogeneity makes this property uniform at relative scale on
a punctured ball about a vertex: cover the unit sphere by finitely
many of the corresponding neighborhoods and rescale.

Let $\rho\ge0$ be a smooth probability kernel supported in the
unit ball.  Away from vertices define
\begin{equation}
 \label{sm:variable-convolution}
 h_2(x)=\int_{\R^3}h_1(x+w(x)z)\rho(z)\,dz,
\end{equation}
where $w$ is smooth and positive there and the integration balls
lie in $\Omega$.  Its derivative is
\begin{equation}
 \label{sm:variable-convolution-derivative}
 Dh_2(x)=\int Dh_1(x+w(x)z)\rho(z)\,dz
 +\int \bigl(Dh_1(x+w(x)z)z\bigr)\otimes Dw(x)\rho(z)\,dz.
\end{equation}
The first integral is in the positive-determinant convex hull just
described.  Taking $w$ and $|Dw|$ sufficiently small locally makes
the second integral a small perturbation, and hence gives
$\det Dh_2>0$.  Formula~\eqref{sm:variable-convolution} is smooth
where $w>0$, as is also seen by writing its smooth variable kernel
in the integration variable $y=x+w(x)z$.

On a small punctured ball at $v$ choose
$w(x)=c_v|x-v|$, with $c_v>0$ sufficiently small.  This is
consistent with the relative-scale convex-hull condition and
makes $h_2-h_1(v)$ homogeneous on a smaller ball.  Set
$h_2(v)=h_1(v)$ there.  The required radius function exists by
a locally finite partition of unity with sufficiently small
positive coefficients.  Near the vertices blend this function
with $c_v|x-v|$ on disjoint annuli.  The annular cutoff terms
are bounded by a constant times $c_v$ if the competing radii
are chosen at that same small scale, so both the radius and
gradient restrictions survive.  This is also an instance of
the local minorant construction in Lemma~\ref{pre:local-tolerances}.

Here is the global error choice.  First take disjoint vertex balls
so small that their volumes, multiplied by the fixed local
gradient bounds to the power $p$, meet a summable error budget.
On these balls the gradients in
Equation~\eqref{sm:variable-convolution-derivative} are bounded independently
of a further decrease of their radii.  On a compact set outside
the balls, the local gradients are bounded and converge at every
piecewise smooth point as $w,|Dw|\to0$.  Dominated convergence
therefore gives an arbitrarily small derivative error on that
compact set.  A locally finite compact covering with summable
budgets makes the total $W^{1,p}$ error as small as prescribed.
The value error follows from local Lipschitz continuity and the
small radius.  These arguments include $p=1$.

\paragraph{The vertices.}
Center a homogeneous vertex ball at the origin and subtract
$h_1(v)$.  Write the resulting map as $F$.  Euler's identity is
$DF(x)x=F(x)$.  Since $DF(x)$ is invertible for $x\ne0$, $F(x)$
never vanishes there.  We can thus write
\begin{equation}
 \label{sm:vertex-form}
 F(r\theta)=rR(\theta)\Phi(\theta),\qquad
 R>0,\quad\theta\in S^2.
\end{equation}
Its positive determinant implies that $\Phi:S^2\to S^2$ is a
smooth orientation-preserving local diffeomorphism.  It is a
covering by compactness, and it is one-sheeted because $S^2$ is
simply connected.

Smale's Theorem gives a smooth isotopy $\Phi_\tau$ from $\Phi$
to a rotation, smooth jointly in $(\tau,\theta)$
\cite[Theorem~6]{Smale1959}; for smooth dependence on the sphere
diffeomorphism see also
\cite[Theorem~1.5]{LiWatts2011}.  Interpolate $R$ through positive
functions $R_\tau$ to a positive constant.  For fixed $\tau$ the
homogeneous map
$F_\tau(r\theta)=rR_\tau(\theta)\Phi_\tau(\theta)$
has determinant
\[
 R_\tau(\theta)^3 J_{S^2}\Phi_\tau(\theta)>0.
\]
The family has bounded first derivatives and a positive
determinant minimum, by compactness.  On an arbitrarily small
ball replace $F$ by $F_{\tau(r)}$, where $\tau$ is 0 near the
outer boundary, 1 near the center, and
$\sup_r|r\tau'(r)|$ is sufficiently small.  Such a transition
is obtained by spreading a fixed cutoff over a sufficiently
long interval of $\log r$.  Its extra derivative term is
bounded by $C|r\tau'(r)|$, so positive determinant persists.
At the center the map is a dilation followed by a rotation.
It is therefore smooth and nonsingular there.

The derivative bound for this last modification depends on the
fixed spherical isotopy but not on the support radius.  Shrinking
that radius gives arbitrarily small summable $W^{1,p}$ costs
over all vertices.  The resulting map $g$ is smooth and has
positive determinant throughout $\Omega$.

Finally, in each of the three operations impose local value
tolerances whose sum is less than
\[
 \min\!\left\{\xi(x),\frac12
 \dist(h(x),\R^3\setminus\Omega')\right\}.
\]
The preceding constructions permit these fine tolerances and a
total Sobolev error less than $\eps$.  Lemma~\ref{sm:proper-degree}
makes $g$ a diffeomorphism onto the original target.  This proves
the Proposition.
\end{proof}

It remains to approximate the locally bi-Lipschitz map by a PL
homeomorphism with strong derivative control.  The next two
subsections establish this reduction.

\subsection{Fine meshes and controlled patch insertion}

Fix a coarse, locally finite Whitney decomposition $\mathcal P$ of
$\Omega$ into closed dyadic cubes with disjoint interiors.  Choose
relatively compact enlarged neighborhoods $U_P$ that are still locally
finite.  Enlarge them a fixed finite number of times when necessary.
Since $H$ is locally bi-Lipschitz, there are constants $K_P\ge1$ such
that
\begin{equation}
 \label{sm:coarse-bilip}
 K_P^{-1}|x-y|\le |H(x)-H(y)|\le K_P|x-y|
 \qquad(x,y\in U_P).
\end{equation}
Enlarge $K_P$ also to bound $|DH|$ almost everywhere on $U_P$ in the
chosen matrix norm.  The constants may incorporate the data from finitely many neighboring
coarse cubes.  All references below to data local to $P$ allow this
finite enlargement, but never an enlargement depending on the final
fine-mesh resolution.

\begin{lemma}[Adapted dyadic meshes]
\label{sm:mesh}
There is an absolute integer $n_0$ with the following property.
Given arbitrary positive local upper bounds on the fine scale, there
is a locally finite dyadic cube decomposition $\mathcal Q$ of $\Omega$
with centers $c_Q$ and side lengths $l_Q$ for which, on writing
\begin{equation}
 \label{sm:patch-box}
 B_Q(s)=c_Q+[-s l_Q,s l_Q]^3,
\end{equation}
the following hold:
\begin{enumerate}[label=\textup{(\roman*)}]
 \item $B_Q(100)\Subset\Omega$, and cubes meeting $B_Q(50)$ have
 side lengths comparable to $l_Q$ by absolute constants.
 \item The normalized cube configurations in $B_Q(6)$ belong to a
 finite list.  There is a conforming affine triangulation of
 $\mathcal Q$ with finitely many normalized nondegenerate shapes
 on these configurations.
 \item The cubes have $n_0$ colors so that the closed boxes $B_Q(6)$
 of a single color are pairwise disjoint.
 \item Attach to $Q$ a coarse index $P(Q)$ containing $c_Q$, using
 a fixed convention on coarse boundaries.  The initial scale
 bounds can be reduced so that all interactions through any fixed
 number of neighboring patch layers are confined to a fixed
 locally finite graph of coarse indices, independent of further
 mesh refinement.  The corresponding patches lie in the
 neighborhoods on which the required bounds in
 Equation~\eqref{sm:coarse-bilip} hold.
\end{enumerate}
The upper bounds on the fine scale remain arbitrarily prescribable.
\end{lemma}

\begin{proof}
Choose a positive size function $s$ on $\Omega$ with sufficiently
small absolute Lipschitz constant, below all the prescribed local
bounds and below $10^{-3}\dist(x,\R^3\setminus\Omega)$.
Lemma~\ref{pre:local-tolerances} supplies such a minorant.  Take the
maximal dyadic cubes for which
\[
 l_Q\le\inf_Qs.
\]
Failure of the condition for the dyadic parent gives
$s(x)\le(2+C\Lip(s))l_Q$ on $Q$.  The small Lipschitz constant
then gives local comparability on $B_Q(50)$.  For example, choosing
it sufficiently small makes every dyadic ratio in this neighborhood
belong to $\{1/2,1,2\}$.  The distance bound puts $B_Q(100)$ inside
$\Omega$.  Maximality gives a decomposition, and the positive lower
bound for $s$ near every compact subset gives local finiteness.

The cube vertices have dyadic coordinates.  Bounded local scale
ratios and bounded normalized distance therefore give only finitely
many normalized configurations.  To triangulate, decompose each
cube face into its square contact facets with adjacent cubes.
Put every contact vertex and hanging subdivision point on the
perimeter of each such square.  Cone these perimeter segments from
the square center, and then cone the resulting triangulation of
each cube boundary from the cube center.  The constructions on
shared facets agree.  They give nondegenerate simplices, and their
normalized shapes range over a finite list.

Local comparability bounds the degree of the graph joining cubes
whose $B_Q(6)$ boxes meet.  A coloring with one more than that
degree proves \textup{(iii)}.  Finally choose the initial local
upper bounds so small relative to the coarse Whitney sizes that
the finitely many required patch layers stay in fixed coarse
neighborhoods.  Since the coarse enlarged cover is locally finite,
its interaction graph is locally finite.  Taking a still smaller
minorant $s$ does not change any of these conclusions.
\end{proof}

The finite lists record actual normalized geometric configurations,
including the positions of faces and vertices.  The same
dyadic-position argument applies in $B_Q(20)$, using the scale
comparability on $B_Q(50)$.  This includes the older protected
boxes needed below.  Indeed, if a box $B_R(t)$ with $t\le2$ meets
$B_Q(5)$, applying scale comparability at the larger cube gives
$l_R/l_Q\in\{1/2,1,2\}$.  Thus
$|c_R-c_Q|_\infty\le5l_Q+2l_R\le9l_Q$, and
$B_R(2)\subset B_Q(13)$.  Shrinking these boxes by one of
finitely many fixed amounts therefore leaves only finitely many
possible intersections and protection margins.

\begin{lemma}[Insertion with protected cores]
\label{sm:insertion}
Consider a patch $Q$ and a finite collection of older protected
boxes with the relative sizes and positions provided by
Lemma~\ref{sm:mesh}.  For every older box prescribe a larger
half-shrunk box and a smaller fully-shrunk box, separated by a
fixed positive normalized margin.  Let $G:\Omega\to\Omega'$ be
a homeomorphism, and let $h_Q$ be a PL embedding on a neighborhood
of $B_Q(3)$.  Assume that $h_Q=G$ on the overlaps with the
half-shrunk boxes.

For every $\lambda>0$ there is an $a>0$, depending only on
$\lambda$, the normalized source configuration, and a local
constant $K_P$ in Equation~\eqref{sm:coarse-bilip}, such that
\begin{equation}
 \label{sm:insertion-smallness}
 \sup_{B_Q(5)}|G-H|<a l_Q,
 \qquad \sup_{B_Q(3)}|h_Q-H|<a l_Q
\end{equation}
implies the following conclusion.  There is a domain
homeomorphism $A$ supported in $B_Q(5)$, with
\[
 \sup|A-\Id|<\lambda l_Q,
\]
such that $G\circ A=h_Q$ on a neighborhood of $B_Q(2)$ and
$G\circ A=G$ on the fully-shrunk older boxes.  The tolerances
are uniform over a finite list of source configurations.  They
do not depend on the derivatives of $G$ or of $h_Q$.
\end{lemma}

\begin{proof}
Scale the patch to $l_Q=1$.  First reduce $a$ using $K_P$ so
that $h_Q(B_Q(3))\subset G(\operatorname{int}B_Q(4))$.
Indeed, for $x\in B_Q(3)$ and $z\in\partial B_Q(4)$,
the distance $|H(z)-H(x)|$ is bounded below by $K_P^{-1}$.
The small errors in Equation~\eqref{sm:insertion-smallness} preserve
this separation.  Degree on $\operatorname{int}B_Q(4)$,
comparing $G$ with $H$ on the boundary, shows that $h_Q(x)$
is in its image.  The same statement holds with a small
margin around $B_Q(3)$.

The embedding
\[
 \eta=G^{-1}\circ h_Q
\]
is consequently defined there.  If $z=\eta(x)$, then
\begin{equation}
 \label{sm:inverse-value-bound}
 |z-x|\le K_P|H(z)-H(x)|
 \le K_P\bigl(|H(z)-G(z)|+|h_Q(x)-H(x)|\bigr)
 <2K_Pa.
\end{equation}
Thus its closeness to the inclusion uses no inverse-Lipschitz
bound for $G$.

Take a fixed closed neighborhood $C$ of $B_Q(2)$ inside
$B_Q(3)$, and let $D$ be the union of the fully-shrunk older
boxes intersected with $B_Q(5)$.  The half/full shrink margins
give neighborhoods on which the map $\eta$ near $C$ agrees
with the identity near $D$.  Include also an identity
neighborhood of $\partial B_Q(5)$, separated from $C$.
Choose these neighborhoods once for each normalized source
configuration.  By Equation~\eqref{sm:inverse-value-bound}, sufficiently
small $a$ makes the union of $\eta$ and these identity maps
an embedding, as in Equation~\eqref{sm:patched-embedding}.  Apply
Lemma~\ref{sm:local-extension} with $B=B_Q(5)$ and displacement
$\lambda$.  Precomposition by the resulting $A$ installs
$h_Q$ on $C$ and preserves the fully-shrunk boxes.  Rescaling
proves the Lemma.
\end{proof}

We record explicitly how tolerances in repeated insertions are chosen.
For a patch insertion supported in $B_Q(5)$,
\begin{equation}
 \label{sm:value-update}
 |G(A(x))-H(x)|
 \le |G(A(x))-H(A(x))|+K_P|A(x)-x|.
\end{equation}
After division by the local cube size, interacting patches introduce
only the bounded ratios from Lemma~\ref{sm:mesh}.  Hence, for a fixed
number of stages, all tolerances can be assigned by backward
recursion: start with the desired final value bounds; choose a small
allowed displacement at the last stage; use
Lemma~\ref{sm:insertion} to obtain its input tolerance; require the
preceding value error to be a small fraction of that tolerance; and
continue backward.  At a coarse index take the minimum over its
finitely many interacting indices at each step.  Only finitely many
steps occur.  Every tolerance remains positive, and none depends on a
PL derivative bound selected later.

\subsection{Finite PL libraries before resolution selection}

The use of finite normalized PL families with exact agreement on
earlier overlaps has a precedent in V\"ais\"al\"a's construction
\cite[Sections~2.7--2.11]{Vaisala1977}.

\begin{proposition}[Strong PL approximation of a locally bi-Lipschitz map]
\label{sm:bilip-pl}
Under the hypotheses of Theorem~\ref{sm:smoothing}, there is a
locally finite PL homeomorphism $G:\Omega\to\Omega'$ satisfying
\[
 \norm{G-H}_{W^{1,p}(\Omega)}<\eps,
 \qquad |G(x)-H(x)|<\xi(x)\quad(x\in\Omega).
\]
\end{proposition}

\begin{proof}
Fix the coarse neighborhoods, local bounds $K_P$, and finite
mesh configurations of Lemma~\ref{sm:mesh}.  We use two rounds of
$n_0$ insertion stages.  The first installs a single piecewise
affine interpolant on the cubes where $H$ is nearly affine; its
gradients there are close to $DH$.  The second installs PL models
on every cube while preserving protected neighborhoods of the
first-round cubes.  It gives
local derivative bounds everywhere, fixed before the remaining
cubes are made small in measure.  For the next steps, let
$\mathcal Q$ be any admissible mesh.  We choose the tolerances
and model libraries uniformly over all such meshes; the mesh
used for the final map will be selected only after the libraries
and their derivative bounds have been fixed.

Write $S=2n_0$ for the number of stages.  A newly installed patch
has protected radius 2.  At each subsequent stage decrease
the radius of every older protected patch by
\begin{equation}
 \label{sm:shrink-size}
 d_*=(4n_0+4)^{-1}
\end{equation}
in its own $B_Q$ coordinates.  Its final radius is greater
than $3/2$, and in particular its closed original cube $Q$
remains in the interior of its protected core.  At an
insertion we use half of the current prescribed decrease
for agreement of the new model and the old map, and the
full decrease for the sets fixed by the ambient change.
These are precisely the margins in
Lemma~\ref{sm:insertion}.

Choose all insertion and displacement tolerances by the
backward procedure following Equation~\eqref{sm:value-update}.
They can be chosen to give an arbitrarily small final
normalized value error on every patch.  If $a_{s,P}$ is
the accuracy required of a new model at stage $s$, require
the current $G-H$ bound before that stage to be less than
$a_{s,P}/8$, also on its interacting neighborhoods.  This
additional fraction is imposed during the same backward
recursion.  All these numbers are now fixed.

\paragraph{Quantized affine interpolation and good patches.}
Choose positive thresholds $\delta_P$ so small that errors
$C\delta_P$ suffice at every first-round insertion involving
that coarse neighborhood, and so that
\begin{equation}
 \label{sm:good-error-budget}
 \sum_{P\in\mathcal P}(C\delta_P)^p |U_P|<\eps^p/8.
\end{equation}
Here and below a local tolerance is reduced using finitely many
neighboring indices when necessary.  We also require
$C\delta_P<(2K_P)^{-1}$, where $C$ is the interpolation
constant for the finite mesh shapes.

We quantize the interpolation data so that their exact values on
protected overlaps, after normalization, have only finitely many
possibilities.  At each vertex $v$ of the fine triangulation, round $H(v)$
to a target dyadic grid whose spacing is a small dyadic
multiple of the smallest incident cube size.  The dyadic
factor is fixed from the coarse indices of those cubes,
small enough that all the corresponding normalized
rounding errors are at most $\delta_P$.  Interpolate
the rounded values affinely on the conforming
triangulation to obtain a continuous piecewise affine
map $L$.  We do not assert that $L$ is globally injective.

Call a cube $Q$, with $P=P(Q)$, good if there is an affine
map $A_Q$ of bi-Lipschitz constant at most $2K_P$ such that
\begin{equation}
 \label{sm:good-tests}
 \sup_{B_Q(20)}|H-A_Q|\le\delta_P l_Q,
 \qquad
 \left(\frac{1}{|Q|}\int_Q|DH-DA_Q|^p\right)^{1/p}\le\delta_P.
\end{equation}
Changing the fixed factor 20 by a larger absolute factor
would have the same effect.  The rounding and the finite
shape list imply
\begin{equation}
 \label{sm:good-interpolation}
 \begin{aligned}
 |DL-DA_Q|&\le C\delta_P
 &&\text{near }B_Q(3),\\
 \sup_{B_Q(3)}|L-H|&\le C\delta_P l_Q.
 \end{aligned}
\end{equation}
For completeness, subtract $A_Q$ at the vertices of any
such tetrahedron.  The vertex errors are at most
$C\delta_P l_Q$; the inverse matrix of its three edge
vectors has norm at most $C/l_Q$, by finite normalized
shapes.  Multiplication gives the first estimate.
The value estimate follows by affine interpolation
and Equation~\eqref{sm:good-tests}.

The first estimate in Equation~\eqref{sm:good-interpolation} makes
$L$ an embedding on a neighborhood of $B_Q(3)$.
Indeed, $L-A_Q$ is Lipschitz on a slightly larger convex
box with constant at most $C\delta_P$; integrate its
piecewise constant derivatives on segments.  Thus
\[
 |L(x)-L(y)|\ge
 \bigl((2K_P)^{-1}-C\delta_P\bigr)|x-y|>0.
\]
Invariance of domain gives the open embedding assertion.

\paragraph{The first round.}
Start with $G_0=H$.  For stages $1,\ldots,n_0$ process the
colors in order, inserting a patch only if it is good,
and taking its model to be $h_Q=L$.  On every overlap
with an older protected core this agrees with the
already installed map, because both equal the single
global map $L$.  Equation~\eqref{sm:good-interpolation}
and the chosen thresholds supply the input accuracy
for Lemma~\ref{sm:insertion}.  Supports $B_Q(5)$ of one
color are disjoint, so these insertions are simultaneous.
Their local finiteness makes the union a domain
homeomorphism.  The value invariants follow from
Equation~\eqref{sm:value-update}.  At the end of this round the
map agrees with $L$ on every good cube and on a
protected neighborhood of it.

\paragraph{Normalized data and their finiteness.}
For each $Q$ choose $m_Q\in\Z^3$ nearest to
$H(c_Q)/l_Q$ and use normalized coordinates
\begin{equation}
 \label{sm:normalization}
 z=\frac{x-c_Q}{l_Q},\qquad
 \widehat H_Q(z)=\frac{H(c_Q+l_Qz)-l_Qm_Q}{l_Q}.
\end{equation}
The value at $z=0$ is bounded by an absolute constant,
and Equation~\eqref{sm:coarse-bilip} gives a uniform local
Lipschitz bound.  The normalized vertex values of $L$
on $B_Q(4)$ lie in a bounded set on finitely many
dyadic grids.  They therefore take only finitely many
values for each coarse index.  Here it matters that
the translation $l_Qm_Q$ is included in the data:
without it one would not have finiteness of affine
maps, as opposed to finiteness of their gradients.

If $R$ interacts with $Q$, then
\begin{equation}
 \label{sm:relative-origins}
 \frac{l_Rm_R-l_Qm_Q}{l_Q}
\end{equation}
is a bounded dyadic vector with a denominator from a
finite list.  Boundedness follows by comparing
$H(c_R)$ and $H(c_Q)$ using Equation~\eqref{sm:coarse-bilip};
the denominators are fixed by the bounded dyadic
scale ratios.  Consequently the first-round exact
PL data, as viewed in any interacting normalized
patch, have only finitely many possibilities.

\paragraph{Choosing the second-round libraries.}
For stages $n_0+1,\ldots,2n_0$ we insert at every
patch of the current color.  The models are chosen
as follows, before the final fine scale is selected.
For each coarse index, the normalized maps
$\widehat H_Q$ on a fixed larger patch form a
uniformly bounded equicontinuous family.  They
admit finitely many sup-norm bins of any prescribed
positive diameter.  At stage $s$ use diameter less
than $a_{s,P}/8$.

Inductively assume that the already installed normalized PL models
have finite lists.  Let $\widehat h_R$ be an older model in its own
$R$-normalization.  In the normalized coordinates of an interacting
patch $Q$, this same map is
\begin{equation}
 \label{sm:transported-model}
 z\longmapsto
 \frac{l_R}{l_Q}\widehat h_R\!\left(
       \frac{c_Q-c_R}{l_R}+\frac{l_Q}{l_R}z\right)
       +\frac{l_Rm_R-l_Qm_Q}{l_Q}.
\end{equation}
The source similarities and protected-core positions have finitely
many possibilities by Lemma~\ref{sm:mesh} and
Equation~\eqref{sm:shrink-size}; the target translations do as well
by Equation~\eqref{sm:relative-origins}.  Only boundedly many older
cores meet $B_Q(5)$, and their coarse indices lie in a fixed finite
neighborhood.  Record each older model's library label, its
normalized source placement and target translation, and its current
protected radius.  These records determine the model on the whole
protected core and have only finitely many possibilities.  In
particular, their restrictions to the half-shrunk overlaps give a
finite list of exact PL data.  No quantization of the later models'
affine coefficients is needed.  There are consequently only finitely
many combinations of
\begin{equation}
 \label{sm:model-combination}
 \text{source configuration, normalized }H\text{-bin,
 and exact protected-model records}.
\end{equation}

For each combination which can occur, select one
representative normalized embedding $G^*$, close
to a map $H^*$ in that bin by the prescribed
current value tolerance, and realizing the
specified older data.  Take the representative
on the fixed open patch $(-4,4)^3$.  If
$\widehat K_i$ are the normalized half-shrunk
older boxes, the set
\[
 K=[-3,3]^3\cap\bigcup_i\widehat K_i
\]
is compact in that open patch.  It lies
strictly inside the currently protected
cores, so $G^*$ is PL on a neighborhood of
$K$.  Apply Lemma~\ref{sm:relative-pl} to
$G^*$ on the open patch, keeping $K$ fixed.
Obtain a PL embedding $h^*$ there with
uniform error less than $a_{s,P}/8$ on the
required compact patch.

The full protected-model records also ensure agreement on a
neighborhood of $K$ with every actual tuple having the same
combination: both its current map and $G^*$ equal the recorded
models on the protected cores, which contain the half-shrunk
boxes in their interiors.  Relative approximation makes $h^*=G^*$
on a neighborhood of $K$.  Since there are only finitely many
closed half-shrunk boxes, we can choose a neighborhood $N$ of
$[-3,3]^3$ inside $(-4,4)^3$ whose intersection with each such box
lies in this common agreement neighborhood.  Store $h^*|_N$ as the library
entry indexed by the coarse region, stage, and combination in
Equation~\eqref{sm:model-combination}.  Its entire overlap with
each half-shrunk box now has the equality required by
Lemma~\ref{sm:insertion}.

If a different actual tuple has the same
combination in Equation~\eqref{sm:model-combination}, this
single model agrees with all of its prescribed
PL overlap data.  Moreover,
\begin{equation}
 \label{sm:reuse-error}
 \norm{h^*-\widehat H_Q}_{\infty}
 \le \norm{h^*-G^*}_{\infty}
     +\norm{G^*-H^*}_{\infty}
     +\norm{H^*-\widehat H_Q}_{\infty}
 <\frac38 a_{s,P}.
\end{equation}
The three norms are on the fixed patch needed
for insertion.  The actual current map is
already within $a_{s,P}/8$ of its $H$, so
Lemma~\ref{sm:insertion} applies after undoing
the normalization.  This installs $h_Q$ while
preserving the fully-shrunk older cores.

This selection does not require compactness of
the family of possible intermediate maps $G$.
Only the $H$-family is binned; the intermediate
map is used as an existence witness for an
embedding realizing the exact finite overlap
data.  All possible normalized meshes and
inputs obeying the fixed coarse bounds may
be included when deciding which combinations
occur.  At any stage the choices for a coarse
index depend only on libraries from strictly
earlier stages at its finitely many interacting
indices, so they can be made simultaneously
over all coarse indices.  Hence the choices are independent of
the eventual mesh resolution.

Each stored map is PL on a neighborhood of a
fixed compact patch and thus has finitely many
affine pieces there.  It has a finite upper
derivative bound, unchanged by the common source and target
rescaling.  There are finitely many
stored maps for each local coarse combination,
and only finitely many stages.  After the second round every cube
retains its own installed model on a protected neighborhood,
because its protected radius remains greater than $3/2$.
Thus the final derivative on that cube comes from a stored PL
model, not from an ambient transition.  This proves by induction
the finiteness of the data needed at the next stage and gives
constants $M_P<\infty$ such that the final map satisfies
\begin{equation}
 \label{sm:model-derivative-bound}
 |DG|\le M_P\quad\hbox{almost everywhere on }Q,
 \qquad P=P(Q).
\end{equation}
The constants $M_P$ are fixed before reducing
the fine scale.  They may be arbitrarily large;
no efficient dependence on $K_P$ is asserted.

The final map is locally PL.  If a single triangulation is
desired, take common refinements of the
finitely many affine subdivisions meeting
each compact set and triangulate the resulting
polyhedral cells.  Local finiteness gives a
locally finite affine triangulation on
$\Omega$.  The map is still onto $\Omega'$
because every stage was a precomposition by
a domain homeomorphism.  On every good cube
the map still equals $L$, since its first-round
protected neighborhood was preserved throughout.

\paragraph{Selecting the resolution and paying for bad cubes.}
Only now choose the final local upper size
bounds of Lemma~\ref{sm:mesh}.  On a fixed
coarse compact region, almost every $x$ is
both a differentiability point of $H$ and a
Lebesgue point of $DH$.  The derivative at
such a point is invertible, with the upper
and lower bounds inherited from local
bi-Lipschitzness.  The affine map
\[
 A_x(y)=H(x)+DH(x)(y-x)
\]
then passes both tests in
Equation~\eqref{sm:good-tests} on every sufficiently
small cube containing $x$.  Indeed its
enlargement is contained in a ball of radius
$C l_Q$ about $x$, which gives the sup-norm
test by differentiability.  The volume of
that ball is at most a fixed multiple of
$|Q|$, which gives the derivative mean test
by the Lebesgue-point property.  Only
finitely many coarse thresholds occur near
a fixed compact set.

It follows that the measure of the union of
bad cubes with a fixed coarse index $P$
tends to zero as their local upper side
length tends to zero.  This statement is
uniform over the admissible meshes.  One
can see the uniformity by taking the union
of all bad dyadic cubes with that index
and side at most $r$.  These measurable
unions decrease as $r\downarrow0$, and
their intersection is null by the preceding
pointwise argument.  They lie in the fixed
finite-volume neighborhood $U_P$.

Choose positive numbers $b_P$ with
$\sum_Pb_P<\eps^p/8$, and reduce the
local mesh bounds until
\begin{equation}
 \label{sm:bad-error-budget}
 (M_P+K_P)^p
 \left|\bigcup_{\substack{Q\text{ bad}\\P(Q)=P}}Q\right|
 <b_P.
\end{equation}
All bounds can be imposed simultaneously
using Lemma~\ref{sm:mesh}.  The $M_P$ in
this inequality were fixed in
Equation~\eqref{sm:model-derivative-bound}; they
are unaffected by this refinement.

The stage tolerances also bound
$|G-H|/l_Q$ by fixed local constants
on each cube.  Further reduce the
local size bounds so that the actual
value error is below $\xi$ and has
$p$th integral less than the remaining
part of $\eps^p$.  For example require
it to be below a global constant
$\eps/[4(1+|\Omega|)^{1/p}]$ and below
the positive minimum of $\xi/2$ on
the relevant compact neighborhood.
These reductions only help the bad-cube
estimate.  Choose a mesh satisfying all these bounds and carry
out the two rounds with the already selected libraries, obtaining
the final map $G$.

On good cubes, Equation~\eqref{sm:good-tests} and
Equation~\eqref{sm:good-interpolation} give
\[
 \int_Q|DG-DH|^p\le(C\delta_P)^p|Q|.
\]
On bad cubes use
Equation~\eqref{sm:model-derivative-bound} and
Equation~\eqref{sm:coarse-bilip}.  Summing and
using Equation~\eqref{sm:good-error-budget} and
Equation~\eqref{sm:bad-error-budget} makes the
derivative error less than the allocated
part of $\eps^p$.

The map is locally Lipschitz,
its derivative is globally $p$-integrable
by the estimates just proved, and its
values lie in the bounded target.
Thus it belongs to $W^{1,p}(\Omega;\R^3)$
and has the required strong error.
\end{proof}

\begin{proof}[Proof of Theorem~\ref{sm:smoothing}]
Apply Proposition~\ref{sm:bilip-pl} with
Sobolev tolerance $\eps/2$ and value
tolerance $\xi/2$, obtaining a PL
homeomorphism $h:\Omega\to\Omega'$.
Apply Proposition~\ref{sm:pl-smoothing}
to $h$ with the same tolerances.
The triangle inequality proves
Equation~\eqref{sm:smoothing-conclusion}.
Both approximations have exactly the
target $\Omega'$.
\end{proof}
